\documentclass[pdf,color,11pt]{CITnote}
%optional [color] for a color logo

\usepackage[margin=30mm]{geometry}
\usepackage{changepage}

\author{Antonio Perreca and ...}
%\refnum{S2....}
\shorttitle{Esperienza 2}
\title{Esperienza x: Gaussian beam profile measurement}
\date{\today}
\issue{1}



\newcommand{\tcr}{\textcolor{red}}
\newcommand{\tcb}{\textcolor{blue}}
\newcommand{\tcm}{\textcolor{magenta}}
\newcommand{\tcg}{\textcolor{green}}
\newcommand{\tcp}{\textcolor{purple}}

\begin{document}

\maketitle

%\section{la}

	%\addtolength{\oddsidemargin}{.875in}
	%\addtolength{\evensidemargin}{.875in}
	
\begin{adjustwidth*}{1.5cm}{1.5cm}
\begin{center}
\section*{Abstract}
\end{center}

{\bf{
Si descrive sinteticamente lo scopo dell' esperienza, cosa e' stato fatto
e i risultati ottenuti (Un paragrafo di 10/15 righe)
}}\\
\end{adjustwidth*}

\section{Introduzione}
Qui si descrive dall' inizio alla fine senza dettagli, in maniera piu' esaustiva
l' intera esperienza (massimo una pagina). Il totale delle pagine dovra' essere di 4 facciate
e non una riga in piu'. La quinta pagina sara' legata alla domanda facoltativa.

\section{Setup}
Si descrive l' apparato di misura e come sono state effettuate
(Lunghezza a scelta...non piu' di una pagina)

\section{Descrizione dei risultati}
Risultati ottenuti ed elaborazione

\begin{figure}[htbp]
	\centering
		\includegraphics[width=14cm]{./fig/Profile_power.pdf}
	\caption{{\bf Profilo longitudinale del fascio gaussiano.}}
	\label{fig:Prof}
\end{figure}

\section{Conclusione}
Si riassumono i risultati ottenuti con la metodica usata. Molto simile all' abstract
ma piu' concentrati sul finale. Un minimo di discussione dei risultati ed eventuali problematiche 
trovate. (10/15 righe al massimo)

\section{Pagina 5: domanda facoltativa}

\end{document}
